% GENERATED FILE. Edit canonical lessons, then run npm run generate:books.
% canonical-source-hash: fnv1a64:8fa148b870f8d906
% canonical-lessons: GU-C212-dharmasala, GU-C212-bhomiyo, GU-C212-rahdari, GU-C212-surang, GU-C212-musapharkhanu, GU-R212-review-207-212, GU-R212-review-final-212

\chapter{Resting on the way}
\label{ch:gu-resting-on-the-way}
\hlchaptermodality{212}

\begin{chapteropening}
\textbf{By the end of this chapter:} \emph{I can say where travellers rested and who guided them.}

Complete the last lesson of chapter 212: I can say where travellers rested and who guided them.
\end{chapteropening}

\section[\texorpdfstring{\gu{ધર્મશાળા} (dharmaśāḷā)}{ધર્મશાળા (dharmaśāḷā)}]{\gu{ધર્મશાળા} (dharmaśāḷā) — rest house, pilgrims' lodge}

\label{lesson:GU-C212-dharmasala}



\begin{quote}
Before the new one: say the Gujarati for pilgrimage, then the Gujarati for vehicle.
\end{quote}

\subsection*{You'll want to know: \gu{ધર્મશાળા}}
\textbf{\gu{ધર્મશાળા}} — \emph{dharmaśāḷā} — \textquotedblleft{}rest house, pilgrims' lodge\textquotedblright{}.

\begin{quote}
\textbf{\gu{અમે} \gu{રાત્રે} \gu{ધર્મશાળામાં} \gu{રહ્યા}.} — \emph{ame rātre dharmaśāḷāmā̃ rahyā} — \textquotedblleft{}We spent the night at a rest house.\textquotedblright{}
\end{quote}

\begin{cousinweb}
\textbf{\gu{ધર્મશાળા}} (\emph{dharmaśāḷā}) is Sanskrit \emph{dharmaśālā}, \textquotedblleft{}charitable asylum\textquotedblright{}, built from \emph{dhárma}, \textquotedblleft{}duty, charity\textquotedblright{}, and \emph{ś\'{\={a}}lā}, \textquotedblleft{}hall\textquotedblright{}, taken over with Gujarati \emph{-ḷ-}, a semi-tatsama (CDIAL 6760a).
\end{cousinweb}

\begin{grammarlens}[title={What you've built}]
One more word for everyday life at A2.
\end{grammarlens}

\subsection*{Your turn}
\begin{itemize}
\raggedright
  \item \emph{Say it:} \emph{dharmaśāḷā}
  \item \emph{Say it:} \emph{dharmaśāḷā}, once more
  \item \emph{Say it:} \emph{ame rātre dharmaśāḷāmā̃ rahyā} — We spent the night at a rest house.
  \item \emph{Recall:} say the Gujarati for pilgrimage, then the Gujarati for vehicle, then say \emph{dharmaśāḷā} again
\end{itemize}

\begin{tcolorbox}[breakable,colback=teal!4,colframe=teal!35!black,title={Before you move on}]
What does \gu{ધર્મશાળા} mean? (\textquotedblleft{}Rest house, pilgrims' lodge\textquotedblright{}.) Say it once more.
\end{tcolorbox}
\section[\texorpdfstring{\gu{ભોમિયો} (bhomiyo)}{ભોમિયો (bhomiyo)}]{\gu{ભોમિયો} (bhomiyo) — guide, one who knows the way}

\label{lesson:GU-C212-bhomiyo}



\begin{quote}
Before the new one: say the Gujarati for vehicle, then the Gujarati for rest house.
\end{quote}

\subsection*{You'll want to know: \gu{ભોમિયો}}
\textbf{\gu{ભોમિયો}} — \emph{bhomiyo} — \textquotedblleft{}guide, one who knows the way\textquotedblright{}.

\begin{quote}
\textbf{\gu{ભોમિયાએ} \gu{અમને} \gu{રસ્તો} \gu{બતાવ્યો}.} — \emph{bhomiyāe amne rasto batāvyo} — \textquotedblleft{}The guide showed us the way.\textquotedblright{}
\end{quote}

\begin{cousinweb}
\textbf{\gu{ભોમિયો}} (\emph{bhomiyo}) is built on \emph{bhom}, \textquotedblleft{}land\textquotedblright{}, with the agent ending \emph{-iyo}: the one who knows the land. CDIAL lists Gujarati \emph{bhom}, \textquotedblleft{}a piece of land, field\textquotedblright{}, under Sanskrit \emph{bhaumika}, \textquotedblleft{}being in the earth\textquotedblright{}, a derivative of \emph{bhū́mi}, \textquotedblleft{}earth\textquotedblright{} (CDIAL 9642).
\end{cousinweb}

\begin{grammarlens}[title={What you've built}]
One more word for everyday life at A2.
\end{grammarlens}

\subsection*{Your turn}
\begin{itemize}
\raggedright
  \item \emph{Say it:} \emph{bhomiyo}
  \item \emph{Say it:} \emph{bhomiyo}, once more
  \item \emph{Say it:} \emph{bhomiyāe amne rasto batāvyo} — The guide showed us the way.
  \item \emph{Recall:} say the Gujarati for vehicle, then the Gujarati for rest house, then say \emph{bhomiyo} again
\end{itemize}

\begin{tcolorbox}[breakable,colback=teal!4,colframe=teal!35!black,title={Before you move on}]
What does \gu{ભોમિયો} mean? (\textquotedblleft{}Guide, one who knows the way\textquotedblright{}.) Say it once more.
\end{tcolorbox}
\section[\texorpdfstring{\gu{રાહદારી} (rāhdārī)}{રાહદારી (rāhdārī)}]{\gu{રાહદારી} (rāhdārī) — pedestrian, passer-by}

\label{lesson:GU-C212-rahdari}



\begin{quote}
Before the new one: say the Gujarati for rest house, then the Gujarati for guide.
\end{quote}

\subsection*{You'll want to know: \gu{રાહદારી}}
\textbf{\gu{રાહદારી}} — \emph{rāhdārī} — \textquotedblleft{}pedestrian, passer-by\textquotedblright{}.

\begin{quote}
\textbf{\gu{એક} \gu{રાહદારીએ} \gu{મને} \gu{સરનામું} \gu{કહ્યું}.} — \emph{ek rāhdārīe mane sarnāmũ kahyũ} — \textquotedblleft{}A passer-by told me the address.\textquotedblright{}
\end{quote}

\begin{cousinweb}
\textbf{\gu{રાહદારી}} (\emph{rāhdārī}) is built from \emph{rāh}, \textquotedblleft{}road, way\textquotedblright{}, usually traced to Persian \emph{rāh}, \textquotedblleft{}road\textquotedblright{}, the Persian suffix \emph{-dār}, \textquotedblleft{}holder\textquotedblright{}, and the ending \emph{-ī}: one who keeps to the road.
\end{cousinweb}

\begin{grammarlens}[title={What you've built}]
One more word for everyday life at A2.
\end{grammarlens}

\subsection*{Your turn}
\begin{itemize}
\raggedright
  \item \emph{Say it:} \emph{rāhdārī}
  \item \emph{Say it:} \emph{rāhdārī}, once more
  \item \emph{Say it:} \emph{ek rāhdārīe mane sarnāmũ kahyũ} — A passer-by told me the address.
  \item \emph{Recall:} say the Gujarati for rest house, then the Gujarati for guide, then say \emph{rāhdārī} again
\end{itemize}

\begin{tcolorbox}[breakable,colback=teal!4,colframe=teal!35!black,title={Before you move on}]
What does \gu{રાહદારી} mean? (\textquotedblleft{}Pedestrian, passer-by\textquotedblright{}.) Say it once more.
\end{tcolorbox}
\section[\texorpdfstring{\gu{સુરંગ} (surang)}{સુરંગ (surang)}]{\gu{સુરંગ} (surang) — tunnel}

\label{lesson:GU-C212-surang}



\begin{quote}
Before the new one: say the Gujarati for guide, then the Gujarati for pedestrian.
\end{quote}

\subsection*{You'll want to know: \gu{સુરંગ}}
\textbf{\gu{સુરંગ}} — \emph{surang} — \textquotedblleft{}tunnel\textquotedblright{}.

\begin{quote}
\textbf{\gu{ટ્રેન} \gu{લાંબી} \gu{સુરંગમાંથી} \gu{ગઈ}.} — \emph{ṭren lā̃bī surangmā̃thī gaī} — \textquotedblleft{}The train went through a long tunnel.\textquotedblright{}
\end{quote}

\begin{cousinweb}
\textbf{\gu{સુરંગ}} (\emph{surang}) is a tadbhava. CDIAL lists it, glossed \textquotedblleft{}underground passage, mine for blasting\textquotedblright{}, under Sanskrit \emph{suruṅgā}, \textquotedblleft{}subterranean passage, mine\textquotedblright{} (CDIAL 13506).
\end{cousinweb}

\begin{grammarlens}[title={What you've built}]
One more word for everyday life at A2.
\end{grammarlens}

\subsection*{Your turn}
\begin{itemize}
\raggedright
  \item \emph{Say it:} \emph{surang}
  \item \emph{Say it:} \emph{surang}, once more
  \item \emph{Say it:} \emph{ṭren lā̃bī surangmā̃thī gaī} — The train went through a long tunnel.
  \item \emph{Recall:} say the Gujarati for guide, then the Gujarati for pedestrian, then say \emph{surang} again
\end{itemize}

\begin{tcolorbox}[breakable,colback=teal!4,colframe=teal!35!black,title={Before you move on}]
What does \gu{સુરંગ} mean? (\textquotedblleft{}Tunnel\textquotedblright{}.) Say it once more.
\end{tcolorbox}
\section[\texorpdfstring{\gu{મુસાફરખાનું} (musāpharkhānũ)}{મુસાફરખાનું (musāpharkhānũ)}]{\gu{મુસાફરખાનું} (musāpharkhānũ) — waiting room}

\label{lesson:GU-C212-musapharkhanu}



\begin{quote}
Before the new one: say the Gujarati for pedestrian, then the Gujarati for tunnel.
\end{quote}

\subsection*{You'll want to know: \gu{મુસાફરખાનું}}
\textbf{\gu{મુસાફરખાનું}} — \emph{musāpharkhānũ} — \textquotedblleft{}waiting room\textquotedblright{}.

\begin{quote}
\textbf{\gu{મુસાફરખાનું} \gu{ભરેલું} \gu{હતું}.} — \emph{musāpharkhānũ bharelũ hatũ} — \textquotedblleft{}The waiting room was full.\textquotedblright{}
\end{quote}

Built on \textbf{\gu{મુસાફર}}, \emph{musāphar}, 'traveller'.

\begin{cousinweb}
\textbf{\gu{મુસાફરખાનું}} (\emph{musāpharkhānũ}) is built from \emph{musāphar}, \textquotedblleft{}traveller\textquotedblright{}, a loan from Arabic \emph{musāfir} that reached Gujarati through Persian, and \emph{khānũ}, \textquotedblleft{}house, room\textquotedblright{}, from Persian \emph{khāna}. The Arabic root is \emph{s-f-r}, \textquotedblleft{}to journey\textquotedblright{} (Lane).
\end{cousinweb}

\begin{grammarlens}[title={What you've built}]
One more word for everyday life at A2.
\end{grammarlens}

\subsection*{Your turn}
\begin{itemize}
\raggedright
  \item \emph{Say it:} \emph{musāpharkhānũ}
  \item \emph{Say it:} \emph{musāpharkhānũ}, once more
  \item \emph{Say it:} \emph{musāpharkhānũ bharelũ hatũ} — The waiting room was full.
  \item \emph{Recall:} say the Gujarati for pedestrian, then the Gujarati for tunnel, then say \emph{musāpharkhānũ} again
\end{itemize}

\begin{tcolorbox}[breakable,colback=teal!4,colframe=teal!35!black,title={Before you move on}]
What does \gu{મુસાફરખાનું} mean? (\textquotedblleft{}Waiting room\textquotedblright{}.) Say it once more.
\end{tcolorbox}
\section[(dialogue)]{Travel by road and rail: a review}

\label{lesson:GU-R212-review-207-212}



\begin{quote}
Say the Gujarati for car and taxi.
\end{quote}

\subsection*{Your turn}
Say these aloud:

\begin{itemize}
\raggedright
  \item car, taxi, driver, petrol, to hire
  \item main road, pothole, traffic, wheel, level crossing
  \item to break down, to push, to collide, to get stuck, to slide
  \item airport, passport, permit, to load, to lift
  \item boat, to ride, to steer, pilgrimage, vehicle
  \item rest house, guide, pedestrian, tunnel, waiting room
\end{itemize}

\begin{tcolorbox}[breakable,colback=teal!4,colframe=teal!35!black,title={Before you move on}]
Car? (\textbf{\gu{ગાડી}}.) Taxi? (\textbf{\gu{ટેક્સી}}.) Driver? (\textbf{\gu{ડ્રાઈવર}}.) Petrol? (\textbf{\gu{પેટ્રોલ}}.) To hire? (\textbf{\gu{ભાડે} \gu{લેવું}}.) Main road? (\textbf{\gu{સડક}}.) Pothole? (\textbf{\gu{ખાડો}}.) Traffic? (\textbf{\gu{ટ્રાફિક}}.) Wheel? (\textbf{\gu{પૈડું}}.) Level crossing? (\textbf{\gu{ફાટક}}.) To break down? (\textbf{\gu{બગડવું}}.) To push? (\textbf{\gu{ધકેલવું}}.) To collide? (\textbf{\gu{અથડાવું}}.) To get stuck? (\textbf{\gu{ફસાવું}}.) To slide? (\textbf{\gu{સરકવું}}.) Airport? (\textbf{\gu{એરપોર્ટ}}.) Passport? (\textbf{\gu{પાસપોર્ટ}}.) Permit? (\textbf{\gu{પરવાનો}}.) To load? (\textbf{\gu{લાદવું}}.) To lift? (\textbf{\gu{ઊંચકવું}}.) Boat? (\textbf{\gu{હોડી}}.) To ride? (\textbf{\gu{સવારી} \gu{કરવી}}.) To steer? (\textbf{\gu{હંકારવું}}.) Pilgrimage? (\textbf{\gu{યાત્રા}}.) Vehicle? (\textbf{\gu{વાહન}}.) Rest house? (\textbf{\gu{ધર્મશાળા}}.) Guide? (\textbf{\gu{ભોમિયો}}.) Pedestrian? (\textbf{\gu{રાહદારી}}.) Tunnel? (\textbf{\gu{સુરંગ}}.) Waiting room? (\textbf{\gu{મુસાફરખાનું}}.)
\end{tcolorbox}
\section[(dialogue)]{Resting on the way — once more}

\label{lesson:GU-R212-review-final-212}



\begin{quote}
Say the Gujarati for rest house and guide.
\end{quote}

\subsection*{Your turn}
\begin{itemize}
\raggedright
  \item \emph{Say it:} rest house, guide, pedestrian, tunnel, waiting room
\end{itemize}

\begin{tcolorbox}[breakable,colback=teal!4,colframe=teal!35!black,title={Before you move on}]
Rest house? (\textbf{\gu{ધર્મશાળા}}.) Guide? (\textbf{\gu{ભોમિયો}}.) Pedestrian? (\textbf{\gu{રાહદારી}}.) Tunnel? (\textbf{\gu{સુરંગ}}.) Waiting room? (\textbf{\gu{મુસાફરખાનું}}.)
\end{tcolorbox}
